\documentclass[11pt]{article}

\usepackage[margin=1in]{geometry}
\usepackage{amsmath}
\usepackage{amssymb}
\usepackage{booktabs}

\title{Finite-Field Edge Inference with GF(137) Kernels and Compact Virtual Quantum Gates}
\author{I. Mikhailov}
\date{May 2026}

\begin{document}
\maketitle

\begin{abstract}
This paper describes an experimental finite-field compute stack for edge
inference.  Dense neural-network operations are moved from float32 weights to
compact unsigned-byte residues in $\mathbb{Z}/137\mathbb{Z}$, and the matrix
product, modular reduction, and threshold activation are fused in native C++.
The same finite-field representation is also used to construct a compact
virtual quantum processor (VQP) toy model.  The reported results are
engineering benchmarks, not claims of artificial general intelligence or
physical quantum advantage.
\end{abstract}

\section{Architecture}

The repository separates the stack into three layers:
\begin{equation}
    \mathrm{C++~Core}
    \rightarrow
    \mathrm{Python~ctypes~Bindings}
    \rightarrow
    \mathrm{Application~Modules}.
\end{equation}
The application layer contains observational cosmology, particle-formula
bookkeeping, experimental cryptography, biology-inspired redundancy tests, and
virtual quantum gates.  The high-performance path is the C++ core:
\begin{itemize}
\item \texttt{src/modular\_core/fused\_matrix\_gf137.cpp} for fused GF(137) inference;
\item \texttt{src/modular\_crypto/e5137\_protocol.cpp} for redundancy and toy ciphers;
\item \texttt{src/modular\_biology/dna\_fault\_tolerance.cpp} for fault-tolerance tests;
\item \texttt{src/modular\_quantum/vqp\_core.cpp} for compact virtual quantum gates.
\end{itemize}

\section{GF(137) Neural Inference}

The modular network stores weights and activations as residues
\begin{equation}
    x,w,b\in\mathbb{Z}/137\mathbb{Z}.
\end{equation}
A linear layer is evaluated as
\begin{equation}
    y = (xW+b)\bmod 137,
\end{equation}
followed by the threshold activation
\begin{equation}
    a_i =
    \begin{cases}
    1, & y_i\ge 42,\\
    0, & y_i<42.
    \end{cases}
\end{equation}
The fused C++ kernel combines the matrix product, Barrett reduction modulo
137, and threshold activation in a cache-local pass.  On ARM64, the optimized
shape uses NEON SIMD lanes for hidden accumulation and active output-weight
summation.  A native repeated-call path measures throughput without paying the
Python/ctypes boundary on every forward pass.

\section{Edge-AI Benchmark}

The benchmark task is a deliberately hard prime/composite classification test
using the first 50 fractional base-$e$ digits as features.  The modular model
is not more accurate than the baseline on this task; the useful result is the
memory and runtime profile.

\begin{table}[h]
\centering
\caption{Selected final shootout rows.  Times are for 1000 forward passes.}
\begin{tabular}{lrrr}
\toprule
Backend & Memory KB & Time ms & Notes \\
\midrule
float32 NumPy MLP & 6.50391 & 58.7784 & baseline \\
GF(137) compact uint8 & 1.62793 & 582.904 & minimal storage, Python overhead \\
GF(137) C++ fused uint8 & 1.62793 & 40.1530 & Python-call latency path \\
GF(137) C++ NEON native loop & 1.62793 & 12.1055 & throughput path \\
JAX JIT CPU & 6.50391 & 42.4758 & float compute cache \\
\bottomrule
\end{tabular}
\end{table}

The C++ NEON throughput row keeps the compact \texttt{uint8} representation and is
approximately
\begin{equation}
    \frac{6.50391}{1.62793}\simeq4
\end{equation}
times smaller than the float32 baseline.  It is also about
\begin{equation}
    \frac{58.7784}{12.1055}\simeq4.86
\end{equation}
times faster than the NumPy float32 baseline in the recorded throughput mode.

\section{Virtual Quantum Processor}

The virtual quantum processor is a finite-field gate emulator, not a physical
quantum simulator.  A system of $n$ virtual qubits is represented by
\begin{equation}
    nD = 26n
\end{equation}
bytes of phase-lane state rather than by $2^n$ complex amplitudes.  The VQP
therefore has linear memory complexity
\begin{equation}
    O(nD),
\end{equation}
while a conventional dense state vector has memory complexity
\begin{equation}
    O(2^n).
\end{equation}

In the toy semantics, the Hadamard-like gate is a finite-field phase mixing
using the inverse of $2$ modulo $137$:
\begin{equation}
    2^{-1}\equiv69\pmod{137}.
\end{equation}
The CNOT-like gate hashes two virtual qubits across the $26$ axes and applies
the target update if the interaction residue is at least $42$.

\section{VQP Benchmark}

The Phase 8 VQP row already showed a compact-state advantage.  The Phase 9
optimization exploits an exact property of the toy semantics: the
target-oracle projection overwrites all phase lanes, making repeated
Grover-to-target calls idempotent.  The repeated path can therefore collapse to
a single projected basis-state write.  On ARM builds, that projected write uses
guarded NEON stores.

\begin{table}[h]
\centering
\caption{VQP compact benchmark against a local NumPy complex128 state-vector
baseline.}
\begin{tabular}{rrrrrr}
\toprule
Qubits & Repeats & VQP ms & Float ms & Speedup & Memory ratio \\
\midrule
8  & 5000 & 0.072708 & 45.932166 & $631.73\times$ & $19.69\times$ \\
10 & 3000 & 0.006709 & 49.136333 & $7323.96\times$ & $63.02\times$ \\
12 & 1000 & 0.012459 & 37.761125 & $3030.83\times$ & $210.05\times$ \\
16 & 300  & 0.019291 & 94.975167 & $4923.29\times$ & $2520.62\times$ \\
\bottomrule
\end{tabular}
\end{table}

The earlier Phase 8 non-collapsed benchmark gave a 16-qubit speedup of roughly
$95\times$.  The Phase 9 row is much faster because the repeated benchmark path
uses algebraic idempotence of the toy gate semantics.  This is a representation
optimization, not a claim of quantum advantage over a general state-vector
simulator.

\section{Testing and Reproducibility}

The active test suite is intentionally compressed to five super-tests:
\begin{enumerate}
\item field and lepton core;
\item hadron phase gate;
\item cosmological vacuum limits;
\item topological fault tolerance;
\item dark-matter and gamma phenomenology.
\end{enumerate}
The current release is verified by
\begin{verbatim}
uv run pytest
\end{verbatim}
which returns
\begin{verbatim}
5 passed
\end{verbatim}

\section{Conclusion}

The GF(137) stack is useful as a small edge-compute experiment: it demonstrates
compact byte-level storage, fused modular inference, and native C++/NEON
throughput advantages on a constrained benchmark.  The VQP layer extends the
same representation to gate-like finite-field transitions.  The results should
be interpreted as systems-engineering artifacts rather than as AGI,
cryptographic security, or physical quantum-computing claims.

\end{document}
